La fonction $g$ produit de fonctions dérivables sur $\R$ est dérivable et sur cet intervalle :\\
$g'(x) = \e^{-x} - x\e^{-x} = \e^{-x}(1 - x)$, qui est elle-même dérivable sur $\R$ et \\
$g''(x) = -\e^{-x} - (1 - x)\e^{-x} = \e^{-x}(x - 2)$.\\
On a $g''(x) = 0 \iff \e^{-x}(x - 2) \iff x - 2 = 0$, puisque $\e^{-x} \ne 0$, quel que soit $x \in \R$, donc la dérivée seconde ne s'annule qu'en $x = 2$.\\
Comme $g(2) = 2\e^{-2} = \dfrac{2}{\e^2}$ et comme $g''(x)$ a le signe de $(x - 2)$ et donc change de signe (de moins à plus) en 2,\\
 le point A$\left(2~;~\dfrac{2}{\e^2}\right)$ est l'unique point d'inflexion de la courbe $\mathcal{C}_g$.\\
\hfill\textbf{Affirmation vraie}